\section{Threats to Validity}
\label{sec:threats}

\subsection{Single-node measurement}

This is the principal threat and we state it without hedging. All measurements
were taken on one machine with a small number of cores. Absolute wall-clock
times therefore do not transfer to a cluster, and we make no claim that they
do.

What survives the objection is specific. Shuffle volume and spill volume are
determined by the physical plan and the data partitioning, not by the machine:
selective salting with $k=32$ replicates the same probe-side rows on a laptop
as on a hundred nodes. The skew ratio is a within-run ratio of task durations
and cancels most machine-specific scaling. And the model's structure ---
convexity in $k$, the $\sqrt{H/P}$ dependence, the two ceilings --- is
analytic, not measured; the experiments test whether it holds, and the
calibration constant $\gamma$ is expected to differ per platform, which is why
it is a calibrated parameter and not a published number.

What does \emph{not} survive is any claim about network shuffle. A single-node
Spark shuffle writes to local disk and reads back; a cluster shuffle crosses the
network, where the replication cost of salting is materially higher. Our
estimate of $b$ is therefore \emph{conservative}, and the true $k^{*}$ on a
cluster should be \emph{smaller} than we report --- which strengthens rather
than weakens the paper's central practical claim that customary salt
cardinalities are too large. We flag this asymmetry explicitly so that a reader
does not have to infer it.

The parallelism ceiling is the sharpest limitation. With few task slots, $C$
binds in Equation~\eqref{eq:keff} for most configurations, so our measurements
exercise that branch more than the cost-optimum branch. We address this by
reporting the model's projection across $C$ rather than only the measured
point, and by being clear that the projection is a projection.

\subsection{Synthetic data}

Zipf is a model of skew, not skew itself. Real key distributions have
structure --- correlated hot keys, temporal drift, keys hot on one side only in
certain periods --- that a stationary Zipf law does not reproduce. We mitigate
this by using an established distributional form, by sweeping $\theta$ rather
than fixing it, by adopting a schema from a real prognostics
dataset~\cite{saxena2008cmapss}, and by including a multi-hot-key
configuration. The harness accepts skewed TPC-H~\cite{microsoft_skewgen} and
real C-MAPSS data as drop-in replacements. But a study on production key
distributions would be stronger, and we do not have one.

\subsection{Bounded probe-side multiplicity}

Our two-sided workloads cap the hot key's probe-side representation at $\mu$
times the mean rather than letting it follow the fact side's Zipf law. This is
a deliberate bound, justified in Section~\ref{sec:methodology}, but it is a
restriction on generality: production joins do exist in which both sides are
severely concentrated and the output genuinely explodes. For those, the binding
problem is output cardinality rather than straggler relief, and neither
mechanism studied here is the right answer --- the correct response is to
change the query, not to tune the join. We regard that as a different problem,
and we do not claim our results speak to it.

\subsection{Version specificity}

AQE's skew-join rule, its defaults, and the interaction between
\texttt{coalescePartitions} and \texttt{skewJoin} all change between Spark
releases, and the production descendant described by Xue et
al.~\cite{xue2024adaptive} differs from the open-source rule. Our measurements
pin one version, recorded in the artifact. The cost model does not depend on
the version; the decision procedure's Branch 2, which reasons about the trigger
condition, does.

\subsection{Construct validity of the skew ratio}

Maximum over median task duration is a proxy, and it has a confound serious
enough that we no longer use it as a primary metric.

The problem is the denominator. AQE's partition coalescing merges the sixteen
shuffle partitions into three, which changes what ``the median task'' refers to
without changing the straggler at all. Comparing a three-task stage's
max-over-median against a sixteen-task stage's is comparing different
quantities. In our data this inverts the ranking: at the thickest probe side
the AQE arm reports a far lower skew ratio than the unmitigated baseline while
its slowest task is \emph{longer} in absolute terms. Any reader who took the
ratio at face value would conclude the opposite of the truth.

We therefore report three things instead. The \emph{absolute maximum task
duration} in the join stage, which is the critical path and has no denominator
to move. The \emph{straggler index}, maximum over \emph{mean} task duration:
because total task time is roughly conserved under coalescing, the mean is
stable where the median is not. And \emph{task counts} alongside every ratio,
so that a reader can see immediately when two numbers are not comparable. The
max-over-median ratio is retained only for continuity with the practitioner
literature, which universally quotes it.

The residual confounds remain: a near-zero median makes the ratio unstable, and
task-duration variance from scheduling and garbage collection is not skew. Both
are reasons to lead with shuffle-byte and record distributions, which have no
such confound.

\subsection{The dimension side is small enough to broadcast}

Our probe table is well under Spark's default ten-megabyte automatic broadcast
threshold. In a default configuration this join would be planned as a broadcast
hash join, which has no shuffle and therefore no shuffle skew, and neither
mechanism we study would be engaged at all. We disable automatic broadcast
throughout precisely to isolate the shuffle-join behaviour we set out to
measure, and we run the broadcast arm so that its dominance at this scale is
visible in the results rather than hidden by the configuration.

The honest scope statement follows: our results apply to fact--dimension joins
in which the dimension side exceeds the broadcast threshold, which is the
regime where skew mitigation is a live question. A reader whose dimension table
fits in a broadcast should broadcast it and stop reading.

\subsection{Execution order and the noise floor}

An earlier version of this harness executed cells in configuration order, which
placed the unmitigated baseline first in every workload block --- on a cold
JVM, and therefore systematically slow. Because every speedup is measured
against that baseline, the bias inflated all of them.

We detected this with a control the harness generates for free: selective
salting at $k=1$ compiles to exactly the same physical plan as the baseline, so
any difference between the two is measurement noise. The discrepancy reached
17.7\% on the first workload while remaining under 0.3\% on the later ones ---
a clear order effect rather than random variation. Cell order is now
randomised under a fixed seed, a session-level warmup precedes the grid, and
the control is computed and reported for every run. It is included in the
artifact.

We state the consequence rather than burying it: \textbf{no speedup smaller
than the control discrepancy should be read as a result}, in this paper or in
any run of this harness.

\subsection{Statistical power}

Five repetitions per cell, and three for the configuration probes, is few. We
report medians and interquartile ranges rather than means and standard
deviations, and we do not attach $p$-values to five-against-five comparisons,
where the smallest attainable two-sided rank-test $p$ is above the conventional
threshold in any case. Differences of a few percent in wall-clock are within
noise at this sample size and we do not interpret them. The deterministic
counters --- shuffle bytes and records --- carry the claims that need
precision, and they do not vary between repetitions at all.

\subsection{One partition count, one core count}

Every reported cell uses sixteen shuffle partitions and two task slots. The
parallelism ceiling therefore binds the recommendation in essentially every
configuration, which means the cost-optimum branch of
Equation~\eqref{eq:keff} --- the branch the closed form exists to serve --- is
never the operative one in our own data. This is the sharpest reason the
cluster validation described in Section~\ref{sec:conclusion} matters, and it
is why we release the cluster configuration alongside the single-node one.

\subsection{Generalisation beyond inner equi-joins}

We evaluate inner equi-joins on a single key. Outer joins, multi-key joins,
joins whose skewed key is also the aggregation key, and joins feeding a
subsequent shuffle all plausibly behave differently. The model's derivation
does not obviously break for these cases, but we have not tested them, and we
do not claim them.

\subsection{Absence claims}

Section~\ref{sec:background} rests on two negative results: no independent
evaluation of Spark's skew-join rule, and no peer-reviewed treatment of
salting. Negative results in literature search are only as good as the search.
Ours covered dblp, the arXiv full-metadata API, Crossref, and domain-scoped
searches across the major publishers, but Semantic Scholar and OpenAlex were
unavailable, and at least one relevant paper~\cite{phan2022skewjoin} is not
indexed in dblp. We therefore present these as strong evidence of a gap rather
than as proof of one, and we cite the nearest existing work explicitly rather
than claiming an empty field.
